\documentclass[12pt,a4paper]{article}

% ---------------- PAQUETES ----------------
\usepackage[spanish]{babel}
\usepackage[utf8]{inputenc}
\usepackage[T1]{fontenc}
\usepackage{geometry}
\usepackage{hyperref}
\usepackage{graphicx}
\usepackage{listings}
\usepackage{xcolor}
\usepackage{longtable}

\geometry{margin=2.5cm}

% ---------------- CONFIGURACION CODIGO ----------------
\definecolor{codegray}{gray}{0.95}
\definecolor{codegreen}{rgb}{0,0.6,0}
\definecolor{codeblue}{rgb}{0,0,0.8}

\lstset{
    backgroundcolor=\color{codegray},
    basicstyle=\ttfamily\small,
    keywordstyle=\color{codeblue},
    commentstyle=\color{codegreen},
    breaklines=true,
    frame=single
}

% ---------------- DOCUMENTO ----------------
\begin{document}

\begin{titlepage}
\centering
{\Large Universidad Mayor de San Andrés \\}
\vspace{0.5cm}
{\large Facultad de Ciencias Puras y Naturales \\}
\vspace{1cm}
{\Huge \textbf{Guía Maestra de Administración y Scripting en Linux} \\}
\vspace{1.5cm}
\textbf{Estudiante:} Andres Brayan Goyzueta Beltran\\ 
\textbf{Materia:} Programación \\ 
\textbf{Docente:} Brigida Carvajal Blanco \\ 
\textbf{Fecha:"23-2-2026} \today
\vfill
\end{titlepage}

\tableofcontents
\newpage

% ==========================================================
\section{Comandos Básicos}
ls,
cd,
pwd,
mkdir,
cp,
mv,
rm,
touch,
cat,
man,
help.
\subsection{Gestión de Directorios}

\subsubsection{ls}
El comando ls permite listar el contenido de un directorio, mostrando archivos y subdirectorios.

\textbf{Ejemplo 1:}
\begin{lstlisting}
ls
\end{lstlisting}

Muestra los archivos del directorio actual.

\textbf{Ejemplo 2:}
\begin{lstlisting}
ls -l
\end{lstlisting}

Muestra los archivos con información detallada como permisos, propietario y tamaño.

\textbf{Ejemplo 3:}
\begin{lstlisting}
ls -a
\end{lstlisting}

Muestra archivos ocultos.


\subsubsection{cd}
El comando cd permite cambiar de directorio dentro del sistema de archivos.

\textbf{Ejemplo 1:}
\begin{lstlisting}
cd Documentos
\end{lstlisting}

Accede al directorio llamado Documentos.

\textbf{Ejemplo 2:}
\begin{lstlisting}
cd ..
\end{lstlisting}

Retrocede al directorio padre.

\textbf{Ejemplo 3:}
\begin{lstlisting}
cd /
\end{lstlisting}

Se mueve al directorio raíz del sistema.
\subsubsection{pwd}
Muestra la ruta absoluta del directorio actual en el que se encuentra el usuario.

\textbf{Ejemplo 1:}
\begin{lstlisting}
pwd
\end{lstlisting}

\textbf{Ejemplo 2:}
\begin{lstlisting}
pwd > ubicacion.txt
\end{lstlisting}

\subsubsection{mkdir}
Permite crear uno o varios directorios.

\textbf{Ejemplo 1:}
\begin{lstlisting}
mkdir proyecto
\end{lstlisting}

\textbf{Ejemplo 2:}
\begin{lstlisting}
mkdir -p proyectos/2025/info
\end{lstlisting}

\subsubsection{cp}
Copia archivos o directorios.

\textbf{Ejemplo 1:}
\begin{lstlisting}
cp archivo.txt copia.txt
\end{lstlisting}

\textbf{Ejemplo 2:}
\begin{lstlisting}
cp -r carpeta respaldo
\end{lstlisting}

\subsubsection{mv}
Mueve o renombra archivos.

\textbf{Ejemplo 1:}
\begin{lstlisting}
mv archivo.txt nuevo.txt
\end{lstlisting}

\textbf{Ejemplo 2:}
\begin{lstlisting}
mv archivo.txt Documentos/
\end{lstlisting}

\subsubsection{rm}
Elimina archivos o directorios.

\textbf{Ejemplo 1:}
\begin{lstlisting}
rm archivo.txt
\end{lstlisting}

\textbf{Ejemplo 2:}
\begin{lstlisting}
rm -r carpeta
\end{lstlisting}

\subsubsection{touch}
Crea archivos vacíos o actualiza su fecha de modificación.

\textbf{Ejemplo 1:}
\begin{lstlisting}
touch nuevo.txt
\end{lstlisting}

\textbf{Ejemplo 2:}
\begin{lstlisting}
touch archivo1.txt archivo2.txt
\end{lstlisting}

\subsubsection{cat}
Muestra el contenido de un archivo en la terminal.

\textbf{Ejemplo 1:}
\begin{lstlisting}
cat archivo.txt
\end{lstlisting}

\textbf{Ejemplo 2:}
\begin{lstlisting}
cat archivo1.txt archivo2.txt
\end{lstlisting}

\subsubsection{man}
Muestra el manual de un comando.

\textbf{Ejemplo 1:}
\begin{lstlisting}
man ls
\end{lstlisting}

\textbf{Ejemplo 2:}
\begin{lstlisting}
man chmod
\end{lstlisting}

\subsubsection{help}
Muestra ayuda interna de comandos integrados en Bash.

\textbf{Ejemplo 1:}
\begin{lstlisting}
help cd
\end{lstlisting}

\textbf{Ejemplo 2:}
\begin{lstlisting}
help echo
\end{lstlisting}

% ==========================================================
\section{Comandos Avanzados}

\subsection{Permisos y Propietarios}

\subsubsection{chmod}
El comando chmod permite cambiar los permisos de lectura (r), escritura (w) y ejecución (x) de un archivo o directorio.

\textbf{Ejemplo 1:}
\begin{lstlisting}
chmod 755 script.sh
\end{lstlisting}
Asigna permisos rwxr-xr-x al archivo.

\textbf{Ejemplo 2:}
\begin{lstlisting}
chmod +x script.sh
\end{lstlisting}
Agrega permiso de ejecución al archivo.


\subsubsection{chown}
Permite cambiar el propietario de un archivo o directorio.

\textbf{Ejemplo 1:}
\begin{lstlisting}
sudo chown usuario archivo.txt
\end{lstlisting}

\textbf{Ejemplo 2:}
\begin{lstlisting}
sudo chown usuario:grupo archivo.txt
\end{lstlisting}


\subsection{Filtros y Procesamiento de Texto}

\subsubsection{grep}
Busca patrones de texto dentro de archivos.

\textbf{Ejemplo 1:}
\begin{lstlisting}
grep "error" archivo.log
\end{lstlisting}

\textbf{Ejemplo 2:}
\begin{lstlisting}
grep -i "usuario" archivo.txt
\end{lstlisting}


\subsubsection{find}
Busca archivos dentro del sistema de archivos.

\textbf{Ejemplo 1:}
\begin{lstlisting}
find . -name "archivo.txt"
\end{lstlisting}

\textbf{Ejemplo 2:}
\begin{lstlisting}
find /home -type f
\end{lstlisting}


\subsubsection{head}
Muestra las primeras líneas de un archivo.

\textbf{Ejemplo 1:}
\begin{lstlisting}
head archivo.txt
\end{lstlisting}

\textbf{Ejemplo 2:}
\begin{lstlisting}
head -n 5 archivo.txt
\end{lstlisting}


\subsubsection{tail}
Muestra las últimas líneas de un archivo.

\textbf{Ejemplo 1:}
\begin{lstlisting}
tail archivo.txt
\end{lstlisting}

\textbf{Ejemplo 2:}
\begin{lstlisting}
tail -f archivo.log
\end{lstlisting}


\subsubsection{sort}
Ordena líneas de texto.

\textbf{Ejemplo 1:}
\begin{lstlisting}
sort archivo.txt
\end{lstlisting}

\textbf{Ejemplo 2:}
\begin{lstlisting}
sort -r archivo.txt
\end{lstlisting}


\subsubsection{wc}
Cuenta líneas, palabras y caracteres.

\textbf{Ejemplo 1:}
\begin{lstlisting}
wc archivo.txt
\end{lstlisting}

\textbf{Ejemplo 2:}
\begin{lstlisting}
wc -l archivo.txt
\end{lstlisting}


\subsection{Procesos del Sistema}

\subsubsection{top}
Muestra procesos en tiempo real y consumo de recursos.

\textbf{Ejemplo:}
\begin{lstlisting}
top
\end{lstlisting}


\subsubsection{ps}
Muestra procesos en ejecución.

\textbf{Ejemplo 1:}
\begin{lstlisting}
ps
\end{lstlisting}

\textbf{Ejemplo 2:}
\begin{lstlisting}
ps aux
\end{lstlisting}


\subsubsection{kill}
Finaliza un proceso mediante su PID.

\textbf{Ejemplo 1:}
\begin{lstlisting}
kill 1234
\end{lstlisting}

\textbf{Ejemplo 2:}
\begin{lstlisting}
kill -9 1234
\end{lstlisting}

% ==========================================================
\section{Programación Shell}

La programación Shell permite automatizar tareas administrativas mediante scripts ejecutados en la terminal.

\subsection{Estructura Básica}

Todo script Bash inicia con:

\begin{lstlisting}
#!/bin/bash
\end{lstlisting}

Esto indica que el intérprete será Bash.

---

\subsection{Variables}

Las variables permiten almacenar datos temporalmente.

\begin{lstlisting}
#!/bin/bash
nombre="Administrador"
echo "Bienvenido $nombre"
\end{lstlisting}

---

\subsection{Estructura Condicional (if)}

Permite ejecutar código según una condición.

\begin{lstlisting}
#!/bin/bash
if [ -f archivo.txt ]; then
    echo "El archivo existe"
else
    echo "El archivo no existe"
fi
\end{lstlisting}

Este ejemplo verifica si un archivo existe en el sistema.

---

\subsection{Bucle for}

Permite repetir acciones sobre un conjunto de elementos.

\begin{lstlisting}
#!/bin/bash
for archivo in *.txt
do
    echo "Procesando $archivo"
done
\end{lstlisting}

Itera sobre todos los archivos .txt del directorio.

---

\subsection{Bucle while}

Ejecuta un bloque mientras la condición sea verdadera.

\begin{lstlisting}
#!/bin/bash
contador=1
while [ $contador -le 5 ]
do
    echo "Numero $contador"
    contador=$((contador+1))
done
\end{lstlisting}

---

\subsection{Estructura case}

Permite seleccionar opciones según el valor de una variable.

\begin{lstlisting}
#!/bin/bash
echo "Seleccione una opcion:"
echo "1) Ver usuarios"
echo "2) Ver fecha"

read opcion

case $opcion in
1)
    who
    ;;
2)
    date
    ;;
*)
    echo "Opcion invalida"
    ;;
esac
\end{lstlisting}

---

\subsection{Script Administrativo Completo}

Ejemplo de script que realiza respaldo automático de archivos:

\begin{lstlisting}
#!/bin/bash

origen="/home/usuario/documentos"
destino="/home/usuario/respaldo"

if [ -d "$origen" ]; then
    cp -r $origen $destino
    echo "Respaldo realizado correctamente"
else
    echo "El directorio origen no existe"
fi
\end{lstlisting}

Este script verifica si existe un directorio y realiza una copia de seguridad.

% ==========================================================
\section{Errores Comunes del Sistema}

Durante la administración y ejecución de comandos en Linux pueden presentarse diversos errores. A continuación se detallan algunos de los más comunes, su causa y solución.

\subsection{Permission Denied}

\textbf{Mensaje:}
\begin{lstlisting}
bash: ./script.sh: Permission denied
\end{lstlisting}

\textbf{Causa:}  
El usuario no tiene permisos de ejecución o acceso al archivo.

\textbf{Solución:}
\begin{lstlisting}
chmod +x script.sh
\end{lstlisting}
o bien ejecutar con privilegios:
\begin{lstlisting}
sudo ./script.sh
\end{lstlisting}


\subsection{Command Not Found}

\textbf{Mensaje:}
\begin{lstlisting}
bash: comando: command not found
\end{lstlisting}

\textbf{Causa:}  
El comando no está instalado o fue escrito incorrectamente.

\textbf{Solución:}
Verificar ortografía o instalar el paquete correspondiente.


\subsection{No Such File or Directory}

\textbf{Mensaje:}
\begin{lstlisting}
bash: archivo.txt: No such file or directory
\end{lstlisting}

\textbf{Causa:}  
El archivo o directorio no existe en la ruta especificada.

\textbf{Solución:}
Verificar la ruta con:
\begin{lstlisting}
ls
\end{lstlisting}


\subsection{Syntax Error}

\textbf{Mensaje:}
\begin{lstlisting}
syntax error near unexpected token
\end{lstlisting}

\textbf{Causa:}  
Error en la estructura del script, como falta de fi, done o esac.

\textbf{Solución:}
Revisar la sintaxis del script cuidadosamente.


\subsection{File Exists}

\textbf{Mensaje:}
\begin{lstlisting}
mkdir: cannot create directory 'proyecto': File exists
\end{lstlisting}

\textbf{Causa:}  
El directorio ya fue creado anteriormente.

\textbf{Solución:}
Utilizar:
\begin{lstlisting}
mkdir -p proyecto
\end{lstlisting}

% ==========================================================
\section{Tuberías y Redireccionamiento}

Las tuberías permiten conectar la salida estándar de un comando con la entrada de otro comando utilizando el símbolo |. 
Esto permite encadenar procesos y automatizar tareas administrativas.

\subsection{Uso de Pipe ( | )}

\textbf{Ejemplo 1: Filtrar archivos específicos}
\begin{lstlisting}
ls -l | grep ".txt"
\end{lstlisting}
Lista archivos detalladamente y filtra solo los que contienen ".txt".

\textbf{Ejemplo 2: Contar procesos activos}
\begin{lstlisting}
ps aux | wc -l
\end{lstlisting}
Cuenta el número total de procesos activos en el sistema.

\textbf{Ejemplo 3: Buscar usuarios conectados}
\begin{lstlisting}
who | sort
\end{lstlisting}
Muestra usuarios conectados ordenados alfabéticamente.


\subsection{Redireccionamiento de Salida ($>$ y $>>$)}

\textbf{Ejemplo 1: Crear archivo con salida}
\begin{lstlisting}
ls > lista.txt
\end{lstlisting}
Guarda la salida del comando en un archivo nuevo.

\textbf{Ejemplo 2: Agregar información}
\begin{lstlisting}
echo "Nuevo registro" >> lista.txt
\end{lstlisting}
Agrega información sin borrar el contenido anterior.


\subsection{Redireccionamiento de Entrada ($<$)}

\textbf{Ejemplo:}
\begin{lstlisting}
wc -l < lista.txt
\end{lstlisting}
Cuenta las líneas del archivo usando redireccionamiento de entrada.

% ==========================================================
\section{Transparencia en el Uso de Inteligencia Artificial}

Para la elaboración de este manual se utilizó la herramienta ChatGPT como apoyo estructural y de redacción técnica. 
El estudiante comprendió y revisó cada sección antes de su inclusión en el documento final.

\begin{longtable}{|p{3cm}|p{5cm}|p{5cm}|}
\hline
\textbf{IA Utilizada} & \textbf{Prompt Enviado} & \textbf{Parte Generada} \\ \hline

ChatGPT & Generar estructura base en formato LaTeX para manual de Linux & Plantilla inicial del documento \\ \hline

ChatGPT & Generar ejemplos de comandos básicos con explicación & Sección de comandos básicos \\ \hline

ChatGPT & Generar comandos avanzados y ejemplos técnicos & Sección de comandos avanzados \\ \hline

ChatGPT & Explicar estructuras de programación Shell con ejemplos administrativos & Sección de programación Shell \\ \hline

ChatGPT & Listar errores comunes del sistema y su solución & Sección de errores comunes \\ \hline

ChatGPT & Explicar uso de tuberías y redireccionamiento & Sección de tuberías \\ \hline

\end{longtable}

\textbf{Nota:} La herramienta fue utilizada como apoyo académico, manteniendo comprensión total del contenido presentado.

\end{document}